\documentclass[10pt]{article}
\usepackage[a4paper,margin=0.82in]{geometry}
\renewcommand{\rmdefault}{cmr}
\usepackage{amsmath,amssymb,amsthm,mathtools,booktabs,enumitem,xcolor,hyperref}
\hypersetup{colorlinks=true,linkcolor=blue!55!black,citecolor=blue!55!black,urlcolor=blue!55!black}
\setlength{\emergencystretch}{3em}
\setlength{\parindent}{0pt}
\setlength{\parskip}{0.52em}
\setlist{leftmargin=1.55em,itemsep=0.22em,topsep=0.22em}
\sloppy

\newtheorem{theorem}{Theorem}[section]
\newtheorem{proposition}[theorem]{Proposition}
\newtheorem{corollary}[theorem]{Corollary}
\newtheorem{lemma}[theorem]{Lemma}
\theoremstyle{definition}
\newtheorem{definition}[theorem]{Definition}
\newtheorem{remark}[theorem]{Remark}

\newcommand{\Zl}{\mathbf Z_\ell}
\newcommand{\Ql}{\mathbf Q_\ell}
\newcommand{\Qlb}{\overline{\mathbf Q}_\ell}
\newcommand{\Dbc}{\mathcal D^b_c}
\newcommand{\Dctf}{\mathcal D^b_{\mathrm{ctf}}}
\newcommand{\Dbcnorm}{\mathcal D^b_{c,\mathrm{norm}}}
\newcommand{\Dctfnorm}{\mathcal D^b_{\mathrm{ctf},\mathrm{norm}}}
\newcommand{\Dbclat}{\mathcal D^b_{c,\mathrm{lat}}}
\newcommand{\Dbclit}{\mathcal D^b_{c,\mathrm{lit}}}
\newcommand{\Dnorm}{\mathcal D_{\mathrm{norm}}}
\newcommand{\Hom}{\operatorname{Hom}}
\newcommand{\End}{\operatorname{End}}
\newcommand{\Kar}{\operatorname{Kar}}
\newcommand{\Rlim}{R\!\varprojlim}
\newcommand{\Rprod}{R\!\prod}
\newcommand{\Tot}{\operatorname{Tot}}
\newcommand{\Cone}{\operatorname{Cone}}
\newcommand{\Cons}{\mathsf{Cons}}
\newcommand{\Perf}{\mathsf{Perf}}
\newcommand{\depends}[1]{\par\smallskip\noindent\textit{Direct prerequisites:} #1.}

\begin{document}
\begin{center}
{\large\bfseries Adic and Rational Coefficient Categories\par}
\medskip
{\normalsize Normalized towers, derived completion, and rationalization\par}
\end{center}

\begin{abstract}
This note constructs integral and rational $\ell$-adic coefficient
categories using ordinary abelian and triangulated categories.  Finite
reduction is built
from K-flat complexes, derived completion is built from the derived inverse
limit of a tower, and the equivalence between complete complexes and
normalized finite-level systems is decomposed into its reduction, limit,
unit, counit, and conservativity statements.  Rational coefficients are
then obtained by central localization and Karoubi completion, followed by
coherent transition functors for finite extensions of coefficient fields.
\end{abstract}

\tableofcontents

\section{The pinned complete discrete valuation ring}

\begin{definition}[Coefficient datum]\label{def:coefficient-datum}
Fix a prime $\ell$ and a finite extension $E/\Ql$ inside a fixed algebraic
closure.  Let
\[
 A=\mathcal O_E,\qquad \mathfrak m=(\varpi),\qquad
 k=A/\mathfrak m,\qquad A_n=A/\varpi^nA\quad(n\geq1).
\]
Here $A$ is a complete discrete valuation ring, $\varpi$ is a selected
uniformizer, and $k$ is finite of characteristic $\ell$.  The topology on
$A$ is the $\varpi$-adic topology, so that the canonical homomorphism
\[
 A\longrightarrow\varprojlim_{n\geq1}A_n,\qquad
 a\longmapsto(a\bmod\varpi^n)_n
\]
is an isomorphism of topological rings.  For $m\leq n$, write
$\rho_{n,m}:A_n\twoheadrightarrow A_m$ for the quotient map.
\end{definition}

Let $X$ be a scheme.  Fix a small basis $\mathcal B_X$ for
$X_{\mathrm{proet}}$, stable under finite fiber products and admitting
refinements of covering families.  All sheaves below are sheaves on this
fixed site.  This convention makes every category and Hom set locally
small.

\begin{definition}[Continuous coefficient presheaves]\label{def:coefficient-presheaves}
For $U\in\mathcal B_X$, let $|U|$ denote its underlying topological space
and set
\[
 A_X^{\mathrm{pre}}(U)=\operatorname{Cont}(|U|,A),\qquad
 (A_n)_X^{\mathrm{pre}}(U)=\operatorname{Cont}(|U|,A_n),
\]
where $A_n$ is discrete.  Addition, multiplication, zero, and one are
defined pointwise.  If $v:V\to U$, restriction is precomposition,
\[
 v^*:A_X^{\mathrm{pre}}(U)\longrightarrow A_X^{\mathrm{pre}}(V),
 \qquad a\longmapsto a\circ |v|,
\]
and similarly at level $n$.  Let $A_X$ and $(A_n)_X$ denote the associated
sheaves of commutative rings.
\end{definition}

\begin{proposition}[Ring-presheaf identities]\label{prop:ring-presheaf}
For every $v:V\to U$, the restriction $v^*$ is a unital ring homomorphism.
The equalities
\[
 (\operatorname{id}_U)^*=\operatorname{id},\qquad
 (v\circ w)^*=w^*v^*
\]
hold for composable arrows $W\xrightarrow{w}V\xrightarrow{v}U$.
Consequently Definition~\ref{def:coefficient-presheaves} gives presheaves
of commutative rings, and their sheafifications are sheaves of commutative
rings.
\depends{Definitions~\ref{def:coefficient-datum} and
\ref{def:coefficient-presheaves}.}
\end{proposition}

\begin{theorem}[Continuous functions satisfy descent]
\label{thm:continuous-functions-sheaf}
The presheaves $A_X^{\mathrm{pre}}$ and $(A_n)_X^{\mathrm{pre}}$ satisfy the
sheaf equalizer condition on the fixed pro-etale basis.  Thus the canonical
maps to their associated sheaves are isomorphisms, and for every
$U\in\mathcal B_X$ one has literal identities
\[
 A_X(U)=\operatorname{Cont}(|U|,A),\qquad
 (A_n)_X(U)=\operatorname{Cont}(|U|,A_n).
\]
The assertion is compatible with restriction, pointwise ring operations,
and the quotient maps $A\to A_n$.
\depends{Definition~\ref{def:coefficient-presheaves}; effective descent of
continuous maps along covering families of the selected pro-etale basis.}
\end{theorem}

\begin{theorem}[Finite quotient and inverse-limit descriptions]
\label{thm:coefficient-comparison}
The quotient $A\to A_n$ induces a surjective morphism of ring sheaves
\[
 q_n:A_X\longrightarrow(A_n)_X
\]
whose kernel is $\varpi^nA_X$.  Thus
\[
 A_X/\varpi^nA_X\xrightarrow{\sim}(A_n)_X.
\]
The maps $q_n$ assemble to an isomorphism of sheaves of rings
\[
 \lambda_X:A_X\xrightarrow{\sim}
       \varprojlim_{n\geq1}(A_n)_X,
 \qquad a\longmapsto(q_n(a))_n.
\]
The inverse limit is computed objectwise on $\mathcal B_X$.  For $A=\Zl$
the resulting pinned ring sheaf is the standard additive pro-etale
$\ell$-adic coefficient sheaf, and its quotient by $\ell^n$ is the standard
$\mathbf Z/\ell^n\mathbf Z$ coefficient sheaf.
\depends{Theorem~\ref{thm:continuous-functions-sheaf}; completeness and
separatedness of the DVR in Definition~\ref{def:coefficient-datum}; local
lifting of a continuous $A_n$-valued function to an $A$-valued function;
objectwise computation of limits of sheaves on the chosen basis.}
\end{theorem}

\begin{definition}[Module sheaves and derived categories]\label{def:ambient-derived}
For any sheaf of commutative rings $R_X$ on the fixed site, let
$\mathsf{Mod}(R_X)$ be the category of sheaves $M$ of abelian groups
equipped with a morphism
\[
 R_X\otimes_{\mathbf Z}M\longrightarrow M
\]
satisfying the unit and associativity laws; arrows are $R_X$-linear sheaf
maps.

A cochain complex $K=(K^i,d^i)$ has $d^{i+1}d^i=0$.  A chain map is a
quasi-isomorphism if it induces an isomorphism on every cohomology sheaf
$H^i(K)=\ker(d^i)/\operatorname{im}(d^{i-1})$.  Let
\[
 \mathsf K(X,R_X)=
 \mathsf{Hot}\bigl(\mathsf{Cpx}(\mathsf{Mod}(R_X))\bigr)
\]
be the homotopy category, whose arrows are chain maps modulo chain homotopy,
and define the ordinary Verdier localization
\[
 D(X,R_X)=\mathsf K(X,R_X)
       [\{\text{quasi-isomorphisms}\}^{-1}].
\]
Write $D(X,A)=D(X,A_X)$ and $D(X,A_n)=D(X,(A_n)_X)$.  Shifts and
distinguished triangles are induced by shifted complexes and mapping
cones in the homotopy category.
\end{definition}

\begin{theorem}[Abelian and Grothendieck properties]\label{thm:module-grothendieck}
The categories $\mathsf{Mod}(A_X)$ and
$\mathsf{Mod}((A_n)_X)$ are abelian Grothendieck categories.  Kernels and
cokernels are the sheaf kernels and cokernels with their induced module
actions.  Filtered colimits are exact, the categories have generators, and
they have enough injectives.  Hence the unbounded derived categories in
Definition~\ref{def:ambient-derived} exist as ordinary triangulated
categories and admit K-injective models.
\depends{Theorem~\ref{thm:coefficient-comparison}; the general theorem on
modules over a ring object in a Grothendieck topos.}
\end{theorem}

\begin{definition}[The central scalar action]\label{def:central-scalar}
For $a\in A$ and $M\in\mathsf{Mod}(A_X)$, multiplication by the constant
section $a$ is an endomorphism $a_M:M\to M$.  Applied degreewise it gives a
chain endomorphism $a_K:K\to K$, hence an endomorphism in $D(X,A)$.  For
every $u:K\to L$,
\[
 u\circ a_K=a_L\circ u.
\]
Thus $a$ is a central natural endomorphism of the identity functor of
$D(X,A)$ and of every full triangulated subcategory considered below.
\end{definition}

\section{K-flat reduction to finite coefficients}

\begin{definition}[Total tensor complex and K-flatness]
\label{def:kflat}
For complexes $P,Q$ of $A_X$-modules, define
\[
 (P\otimes_{A_X}Q)^r=\bigoplus_{i+j=r}P^i\otimes_{A_X}Q^j,
\]
with differential
\[
 d(p\otimes q)=d_Pp\otimes q+(-1)^i p\otimes d_Qq
 \qquad(p\in P^i).
\]
A complex $P$ is K-flat if $P\otimes_{A_X}C$ is acyclic for every acyclic
complex $C$.  An adapted K-flat resolution of a complex $K$ is a
quasi-isomorphism $p:P\to K$ for which $P$ is termwise flat and K-flat.
\end{definition}

\begin{theorem}[Existence of adapted K-flat sources]\label{thm:kflat-existence}
Every complex $K\in\mathsf{Cpx}(\mathsf{Mod}(A_X))$ admits an adapted
K-flat resolution
\[
 p:P\xrightarrow{\sim}K.
\]
For every quasi-isomorphism $u:K\to L$, the categories of roofs
\[
 K\xleftarrow{\sim}P\longrightarrow L
\]
with K-flat adapted source are cofinal in the corresponding Verdier
denominator categories.  No strict functorial K-flat endofunctor is required
or asserted.
\depends{Theorem~\ref{thm:module-grothendieck}; the K-flat resolution
and adapted-localization theorems for complexes of modules on a small
ringed site.}
\end{theorem}

\begin{theorem}[Derived scalar extension by adapted localization]
\label{thm:derived-scalar-extension}
Let $A_X\to B_X$ be a morphism of sheaves of commutative rings.  On an
adapted K-flat source define
\[
 P\longmapsto P\otimes_{A_X}B_X.
\]
This functor sends quasi-isomorphisms between adapted sources to
quasi-isomorphisms and therefore descends, by the Verdier localization
universal property, to a choice-independent exact triangulated functor
\[
 -\otimes^L_{A_X}B_X:D(X,A_X)\longrightarrow D(X,B_X).
\]
For an adapted resolution $P\xrightarrow{\sim}K$, the value on $K$ is
represented by $P\otimes_{A_X}B_X$.  Successive scalar extensions carry
canonical unit and associativity isomorphisms satisfying the triangle and
pentagon equations.
\depends{Theorem~\ref{thm:kflat-existence}; K-flat tensor preservation of
quasi-isomorphisms; the universal property of Verdier localization.}
\end{theorem}

\begin{definition}[Internal Hom and global sections on a ringed site]
\label{def:ambient-hom-sections}
Let $R$ be a commutative ring and let $R_X$ be a sheaf of commutative
$R$-algebras on the fixed pro-etale site.  For
$M,N\in\mathsf{Mod}(R_X)$, define the internal-Hom sheaf by
\[
 \mathcal Hom_{R_X}(M,N)(U)=
 \Hom_{\mathsf{Mod}(R_X|_U)}(M|_U,N|_U).
\]
Restriction of morphisms makes the right side a sheaf on the site.  For
cochain complexes $P,I$, define the total internal-Hom complex by
\[
 \mathcal Hom_{R_X}^{\bullet}(P,I)^r
 =\prod_{a\in\mathbf Z}
   \mathcal Hom_{R_X}(P^a,I^{a+r}),
 \qquad
 d(f)=d_I\circ f-(-1)^r f\circ d_P.
\]
A complex $I$ is K-injective if the ordinary mapping complex
$\operatorname{Hom}^{\bullet}_{\mathsf{Cpx}(R_X)}(C,I)$ is acyclic for
every acyclic complex $C$.

The global-sections functor is
\[
 \Gamma(X,-):\mathsf{Mod}(R_X)\longrightarrow\mathsf{Mod}(R),
 \qquad M\longmapsto M(X),
\]
where the $R$-action is induced by $R\to\Gamma(X,R_X)$.  Let $D(R)$ be
the Verdier localization of the homotopy category of complexes of
$R$-modules at quasi-isomorphisms.
\end{definition}

\begin{theorem}[Derived internal Hom and derived global sections]
\label{thm:ambient-hom-sections}
Assume that $\mathsf{Mod}(R_X)$ is Grothendieck abelian and has the adapted
K-flat and K-injective resolution properties.  There are ordinary derived
functors
\[
 R\mathcal Hom_{R_X}^{\mathrm{amb}}:
 D(X,R_X)^{\mathrm{op}}\times D(X,R_X)
 \longrightarrow D(X,R_X)
\]
and
\[
 R\Gamma(X,-):D(X,R_X)\longrightarrow D(R).
\]
If $P\xrightarrow{\sim}K$ is adapted K-flat and
$L\xrightarrow{\sim}I$ is K-injective, then
$R\mathcal Hom_{R_X}^{\mathrm{amb}}(K,L)$ is represented by
$\mathcal Hom_{R_X}^{\bullet}(P,I)$, and $R\Gamma(X,L)$ is represented by
$\Gamma(X,I)$.  The resulting objects and morphisms are independent of
the adapted resolutions in the derived categories.

For $K,L,M\in D(X,R_X)$ there is a natural tensor--Hom adjunction
\[
 \Hom_{D(X,R_X)}
 \bigl(M,R\mathcal Hom_{R_X}^{\mathrm{amb}}(K,L)\bigr)
 \xrightarrow{\sim}
 \Hom_{D(X,R_X)}(M\otimes^L_{R_X}K,L).
\]
For every $r\in\mathbf Z$ there is a natural isomorphism
\[
 \Hom_{D(X,R_X)}(R_X,K[r])
 \xrightarrow{\sim}H^rR\Gamma(X,K).
\]
For finite coefficient sheaves and literal coefficient sheaves, the
superscript $\mathrm{amb}$ is omitted when no completed internal-Hom
construction is being compared.
\depends{Definitions~\ref{def:ambient-derived}, \ref{def:kflat}, and
\ref{def:ambient-hom-sections}; Theorems~\ref{thm:module-grothendieck}
and \ref{thm:kflat-existence}; the closed tensor--Hom adjunction for
complexes of sheaves.}
\end{theorem}

\begin{definition}[The two-term quotient resolution]\label{def:quotient-resolution}
For $n\geq1$, let $L_n$ be the complex of $A_X$-modules
\[
 L_n=\bigl[A_X\xrightarrow{\ \varpi^n\ }A_X\bigr],
\]
with the first copy in degree $-1$ and the second in degree $0$.  The
augmentation $L_n\to(A_n)_X[0]$ is zero in degree $-1$ and is $q_n$ in
degree $0$.
\end{definition}

\begin{proposition}[Finite quotient resolution]\label{prop:quotient-resolution}
Multiplication by $\varpi^n$ on $A_X$ is injective, the augmentation
\[
 L_n\longrightarrow(A_n)_X[0]
\]
is a quasi-isomorphism, and $L_n$ is K-flat.  Consequently $(A_n)_X$ has
Tor-dimension contained in $[-1,0]$ over $A_X$.
\depends{Theorem~\ref{thm:coefficient-comparison};
Definition~\ref{def:quotient-resolution}; the two-term filtration of
$L_n\otimes C$.}
\end{proposition}

\begin{definition}[Ambient finite reduction]\label{def:finite-reduction}
Apply Theorem~\ref{thm:derived-scalar-extension} to the quotient
$A_X\to(A_n)_X$ and define
\[
 \operatorname{red}_n:D(X,A)\longrightarrow D(X,A_n),
 \qquad
 K\longmapsto K\otimes^L_{A_X}(A_n)_X.
\]
If $P\xrightarrow{\sim}K$ is any adapted K-flat resolution, the value is
represented by $P\otimes_{A_X}(A_n)_X$.  After restriction of scalars to
$A_X$, it is canonically isomorphic in $D(X,A)$ to
$P\otimes_{A_X}L_n$; this last complex is used only as an $A_X$-model.
\end{definition}

\begin{theorem}[Well-defined finite reduction]\label{thm:finite-reduction}
Definition~\ref{def:finite-reduction} gives an exact triangulated functor,
independent up to a unique coherent natural isomorphism of every adapted
K-flat resolution.  It sends quasi-isomorphisms to isomorphisms, commutes
with shifts, and sends a cone triangle to the cone triangle of the reduced
map.
\depends{Theorem~\ref{thm:derived-scalar-extension};
Proposition~\ref{prop:quotient-resolution}; the cone criterion for
quasi-isomorphisms.}
\end{theorem}

\begin{definition}[Transitions among reductions]\label{def:reduction-transitions}
For $m\leq n$, derived extension of scalars along
$\rho_{n,m}:(A_n)_X\to(A_m)_X$ is
\[
 \operatorname{red}_{n,m}:D(X,A_n)\longrightarrow D(X,A_m),
 \qquad C\longmapsto C\otimes^L_{(A_n)_X}(A_m)_X.
\]
Associativity of derived tensor product gives
\[
 \alpha_{n,m}(K):
 \operatorname{red}_{n,m}(\operatorname{red}_nK)
       \xrightarrow{\sim}\operatorname{red}_mK.
\]
For $r\leq m\leq n$, the two composites from
$\operatorname{red}_{m,r}\operatorname{red}_{n,m}\operatorname{red}_nK$
to $\operatorname{red}_rK$ are required to agree.
\end{definition}

\begin{theorem}[Reduction coherence]\label{thm:reduction-coherence}
The maps $\alpha_{n,m}$ are natural isomorphisms.  Together with the
identity comparison at $m=n$, they satisfy the unit and pentagon equations.
Thus all finite reductions form a coherent inverse system of exact
triangulated functors.
\depends{Theorem~\ref{thm:finite-reduction}; associativity and unit
coherence for derived tensor product.}
\end{theorem}

\section{Finite-level constructibility and uniform systems}

\begin{definition}[Finite-level lisse and constructible objects]
\label{def:finite-constructible}
An $(A_n)_X$-module sheaf is \emph{lisse} if it is locally isomorphic on
$X_{\mathrm{proet}}$ to the constant sheaf associated with a finite
$A_n$-module.  It is \emph{finite locally free} if the module may be chosen
finite free.  It is \emph{constructible} if every affine open of $X$ has a
finite partition by constructible locally closed subschemes on which the
restriction is lisse.

Let $\Dbc(X,A_n)$ be the full subcategory of $D(X,A_n)$ consisting of
complexes whose cohomology sheaves are constructible and vanish outside a
finite interval.  Such a complex has Tor-amplitude in $[a,b]$ if
\[
 H^i(C\otimes^L_{(A_n)_X}M)=0
 \quad\text{for every module sheaf $M$ and every $i\notin[a,b]$}.
\]
The full subcategory having some finite Tor-amplitude is
$\Dctf(X,A_n)$.
\end{definition}

\begin{theorem}[Finite-level closure]\label{thm:finite-closure}
Assume that $X$ is topologically noetherian.  Constructible
$(A_n)_X$-module sheaves form an abelian subcategory closed under extensions
and direct summands.  Hence $\Dbc(X,A_n)$ and $\Dctf(X,A_n)$ are full
triangulated subcategories of $D(X,A_n)$.  Finite scalar extension and
derived tensor with a perfect coefficient complex preserve
$\Dctf(X,A_n)$.
\depends{Definition~\ref{def:finite-constructible}; finite stratification
refinement; finite-level constructibility for torsion etale sheaves and the
etale--pro-etale comparison.}
\end{theorem}

\begin{definition}[The abelian category of coefficient systems]
\label{def:system-abelian}
Let $\mathsf{Sh}(X)$ denote the category of sheaves of abelian groups on the
fixed site.  The quotient maps make
\[
 A_\bullet=((A_n)_X)_{n\geq1}
\]
a commutative ring object of
$\mathsf{Sh}(X)^{\mathbf N_{\geq1}^{\mathrm{op}}}$.  Define
$\mathsf{Mod}(A_\bullet)$ to be the category of modules over this ring
object.  Explicitly, an $A_\bullet$-module sheaf is a sequence
\[
 M_\bullet=(M_n,t_{n+1,n})_{n\geq1},
\]
where $M_n\in\mathsf{Mod}((A_n)_X)$ and
\[
 t_{n+1,n}:M_{n+1}\longrightarrow M_n
\]
is $(A_{n+1})_X$-linear after restriction of scalars along
$(A_{n+1})_X\to(A_n)_X$.  A morphism $u_\bullet:M_\bullet\to N_\bullet$ is
a family $u_n:M_n\to N_n$ satisfying
\[
 u_n t^M_{n+1,n}=t^N_{n+1,n}u_{n+1}.
\]
Let
\[
 \mathsf K(X,A_\bullet)=
 \mathsf{Hot}\bigl(\mathsf{Cpx}(\mathsf{Mod}(A_\bullet))\bigr)
\]
and define the ordinary Verdier localization
\[
 D(X,A_\bullet)=
 \mathsf K(X,A_\bullet)
 [\{\text{levelwise quasi-isomorphisms}\}^{-1}].
\]
Evaluation at $n$ is denoted
$\operatorname{ev}_n:D(X,A_\bullet)\to D(X,A_n)$.
\end{definition}

\begin{theorem}[Abelian system category]\label{thm:system-abelian}
The category $\mathsf{Mod}(A_\bullet)$ is abelian and Grothendieck.  Limits,
colimits, kernels, and cokernels are computed levelwise with induced
transition maps.  Every $\operatorname{ev}_n$ is exact, hence induces the
displayed exact triangulated evaluation functor.
\depends{Theorem~\ref{thm:module-grothendieck};
Definition~\ref{def:system-abelian}; the Grothendieck theorem for modules
over a ring object in a countable diagram category.}
\end{theorem}

\begin{definition}[Normalized derived system]\label{def:normalized-system}
For $K_\bullet\in D(X,A_\bullet)$ and $m\leq n$, the transition maps define
\[
 \theta_{n,m}(K_\bullet):
 \operatorname{ev}_n(K_\bullet)
   \otimes^L_{(A_n)_X}(A_m)_X
 \longrightarrow\operatorname{ev}_m(K_\bullet).
\]
The system is \emph{normalized} if every $\theta_{n,m}$ is an isomorphism
and, for $r\leq m\leq n$, the direct comparison $\theta_{n,r}$ equals the
composite through level $m$ after inserting the tensor associator.  Let
$\Dnorm(X,A_\bullet)$ be the full subcategory of normalized systems.
\end{definition}

\begin{theorem}[Triangulated normalization]\label{thm:normalized-triangulated}
The category $\Dnorm(X,A_\bullet)$ is a full triangulated subcategory of
$D(X,A_\bullet)$.  A morphism of normalized systems is an isomorphism if
and only if each evaluated morphism is an isomorphism.
\depends{Theorem~\ref{thm:system-abelian}; exactness of evaluation and
derived scalar extension; the two-out-of-three property for a morphism of
distinguished triangles.}
\end{theorem}

\begin{definition}[Uniform lissity and constructibility]
\label{def:uniform-constructible}
A normalized system $K_\bullet$ is \emph{uniformly lisse} if there is a
pro-etale covering $\{U_i\to X\}$, one finite interval $[a,b]$, and for
each $i$ a bounded complex $P_i$ of finite projective $A$-modules supported
in $[a,b]$, together with isomorphisms, compatible in $n$,
\[
 K_n|_{U_i}\xrightarrow{\sim}
 \underline{P_i\otimes_A^LA_n}_{\,U_i}
 \quad\text{in }D(U_i,A_n).
\]
It is \emph{uniformly constructible} if there is one finite constructible
stratification
\[
 X=\coprod_{\alpha=1}^r X_\alpha
\]
such that its restriction to every $X_\alpha$ is uniformly lisse.  It is
\emph{uniformly finite-Tor} if the complexes $P_i$ may be chosen in one
common Tor-amplitude interval.

Write $\Dbcnorm(X,A_\bullet)$ for the uniformly constructible
normalized systems and $\Dctfnorm(X,A_\bullet)$ for the
uniformly finite-Tor ones.
\end{definition}

\begin{proposition}[Uniform Tor amplitude over a DVR]
\label{prop:uniform-tor-dvr}
Every uniformly lisse system of
Definition~\ref{def:uniform-constructible} has the common Tor-amplitude
interval $[a,b]$ exhibited there.  The same assertion holds stratum by
stratum for a uniformly constructible system.  Consequently, for the
pinned DVR coefficients used in this note,
\[
 \Dbcnorm(X,A_\bullet)=\Dctfnorm(X,A_\bullet).
\]
The two notations are retained in order to record whether a later theorem
uses only constructibility or explicitly uses finite Tor amplitude.
\depends{Definition~\ref{def:uniform-constructible}; flatness of finite
projective $A$-modules and stability of Tor amplitude under derived
reduction $A\to A_n$.}
\end{proposition}

\begin{theorem}[Uniform closure and common refinements]
\label{thm:uniform-closure}
Assume that $X$ is topologically noetherian.  Any two finite constructible
stratifications admit the finite common refinement formed by their nonempty
intersections.  Uniformly lisse systems are closed under shifts, cones,
finite sums, and direct summands.  Consequently
\[
 \Dbcnorm(X,A_\bullet),\qquad
 \Dctfnorm(X,A_\bullet)
\]
are full triangulated subcategories of
$\Dnorm(X,A_\bullet)$.
\depends{Proposition~\ref{prop:uniform-tor-dvr};
Theorems~\ref{thm:finite-closure} and
\ref{thm:normalized-triangulated}; Definition~\ref{def:uniform-constructible};
finite common refinement of constructible stratifications.}
\end{theorem}

\begin{definition}[The reduction-system functor]\label{def:reduction-system}
On complexes, let
\[
 \mathcal T_A(P)=
 \bigl(P\otimes_{A_X}(A_n)_X\bigr)_{n\geq1},
\]
with strict transitions induced by $(A_{n+1})_X\to(A_n)_X$.  On adapted
K-flat sources this construction sends quasi-isomorphisms to levelwise
quasi-isomorphisms.  Its adapted left-derived functor is
\[
 \operatorname{Red}_A:D(X,A)\longrightarrow D(X,A_\bullet),
 \qquad
 \operatorname{ev}_n\operatorname{Red}_A(K)=\operatorname{red}_n(K).
\]
\end{definition}

\begin{proposition}[Reductions are normalized]\label{prop:reductions-normalized}
For every $K\in D(X,A)$, the system $\operatorname{Red}_A(K)$ is normalized.
Its normalization isomorphism at $m\leq n$ is
$\alpha_{n,m}(K)$ from Definition~\ref{def:reduction-transitions}.  Thus
\[
 \operatorname{Red}_A:D(X,A)\longrightarrow
      \Dnorm(X,A_\bullet)
\]
is exact triangulated.
\depends{Theorems~\ref{thm:derived-scalar-extension},
\ref{thm:finite-reduction}, and
\ref{thm:reduction-coherence}; Definition~\ref{def:reduction-system}.}
\end{proposition}

\section{Derived inverse limits and completion}

\begin{definition}[Tower categories and restriction of scalars]
\label{def:tower-category}
For every sheaf of commutative rings $R_X$ on the fixed site, let
\[
 \mathsf{Tow}_{R_X}(X)=
 \operatorname{Fun}(\mathbf N_{\geq1}^{\mathrm{op}},
                    \mathsf{Mod}(R_X)).
\]
An object is a tower $M_1\leftarrow M_2\leftarrow\cdots$ of $R_X$-module
sheaves.  Its inverse limit is the kernel
\[
 \varprojlim_nM_n\longrightarrow\prod_{n\geq1}M_n
 \xrightarrow{\,1-s\,}\prod_{n\geq1}M_n,
 \qquad s((x_n)_n)=(t_{n+1,n}(x_{n+1}))_n.
\]
Define
\[
 \mathsf K(\mathsf{Tow}_{R_X}(X))=
 \mathsf{Hot}\bigl(\mathsf{Cpx}(\mathsf{Tow}_{R_X}(X))\bigr)
\]
and
\[
 D(\mathsf{Tow}_{R_X}(X))=
 \mathsf K(\mathsf{Tow}_{R_X}(X))
 [\{\text{quasi-isomorphisms}\}^{-1}],
\]
where a quasi-isomorphism is detected by every level evaluation.

For $R_X=A_X$, write $\mathsf{Tow}_A(X)$.  The quotient $q_n:A_X\to(A_n)_X$
gives an exact restriction functor
\[
 \operatorname{Res}_{A/A_n}:D(X,A_n)\longrightarrow D(X,A).
\]
Restriction of scalars defines an exact functor
\[
 \operatorname{Res}_A:\mathsf{Mod}(A_\bullet)\longrightarrow
 \mathsf{Tow}_A(X).
\]
The exact restriction functor induces
\[
 \mathbf D\operatorname{Res}_A:D(X,A_\bullet)\longrightarrow
 D(\mathsf{Tow}_A(X)).
\]
\end{definition}

\begin{theorem}[Tower category and derived inverse limit]
\label{thm:derived-limit-existence}
For every $R_X$, the category $\mathsf{Tow}_{R_X}(X)$ is Grothendieck
abelian.  The left exact functor
$\varprojlim:\mathsf{Tow}_{R_X}(X)\to\mathsf{Mod}(R_X)$ therefore has a
total right-derived functor
\[
 \Rlim_{R_X}:D(\mathsf{Tow}_{R_X}(X))\longrightarrow D(X,R_X),
\]
computed by a functorial K-injective resolution in the tower category.
This construction is independent, up to a canonical natural isomorphism,
of the selected K-injective resolution.
\depends{Theorem~\ref{thm:module-grothendieck};
Definition~\ref{def:tower-category}; the unbounded K-injective resolution
theorem.}
\end{theorem}

\begin{theorem}[Levelwise derived reduction of a tower]
\label{thm:tower-reduction}
For every $m\geq1$, derived scalar extension along
$A_X\to(A_m)_X$ defines an exact functor
\[
 \mathsf T_m:D(\mathsf{Tow}_A(X))\longrightarrow
 D(\mathsf{Tow}_{(A_m)_X}(X))
\]
and natural isomorphisms
\[
 \operatorname{ev}_n(\mathsf T_mC_\bullet)
 \xrightarrow{\sim}
 \operatorname{ev}_n(C_\bullet)\otimes^L_{A_X}(A_m)_X.
\]
The induced transition maps are obtained by tensoring the transition maps
of $C_\bullet$, and the isomorphisms are compatible with all evaluations.
\depends{Theorem~\ref{thm:derived-scalar-extension}; the adapted K-flat
localization theorem in the Grothendieck category of towers; exactness of
level evaluation.}
\end{theorem}

\begin{theorem}[Roos--Milnor triangle]\label{thm:milnor-triangle}
Let $C_\bullet\in D(\mathsf{Tow}_{R_X}(X))$ and put
$C_n=\operatorname{ev}_n(C_\bullet)$.  Let
\[
 s:\Rprod_{n\geq1}C_n\longrightarrow\Rprod_{n\geq1}C_n
\]
be induced by the transition maps.  There is a functorial distinguished
triangle
\[
 \Rlim_{R_X} C_\bullet\longrightarrow\Rprod_{n\geq1}C_n
 \xrightarrow{\,1-s\,}\Rprod_{n\geq1}C_n
 \longrightarrow(\Rlim_{R_X}C_\bullet)[1].
\]
If countable products are exact on the full subcategory containing all
terms of the chosen tower, $\Rprod$ may be replaced by the ordinary product.
In that case there are functorial Milnor exact sequences
\[
 0\longrightarrow R^1\!\varprojlim_n H^{i-1}(C_n)
 \longrightarrow H^i(\Rlim_{R_X}C_\bullet)
 \longrightarrow\varprojlim_nH^i(C_n)\longrightarrow0.
\]
\depends{Theorem~\ref{thm:derived-limit-existence}; the two-term Roos
resolution of inverse limit on a countable tower.}
\end{theorem}

\begin{theorem}[Perfect tensor exchange with products and limits]
\label{thm:perfect-tensor-limit}
Let $P$ be a perfect $R_X$-complex, represented locally by a bounded complex
of finite projective $R_X$-modules, and let
$C_\bullet\in D(\mathsf{Tow}_{R_X}(X))$.  The canonical maps
\[
 \left(\Rprod_n C_n\right)\otimes^L_{R_X}P
 \xrightarrow{\sim}
 \Rprod_n(C_n\otimes^L_{R_X}P)
\]
and
\[
 \left(\Rlim_{R_X}C_\bullet\right)\otimes^L_{R_X}P
 \xrightarrow{\sim}
 \Rlim_{R_X}\bigl(C_\bullet\otimes^L_{R_X}P\bigr)
\]
are isomorphisms, natural in $C_\bullet$ and $P$.  Here the tower on the
right is the levelwise derived tensor tower.
\depends{Theorem~\ref{thm:milnor-triangle}; finite projective tensor
commutes with products; induction on the length of a perfect complex and
local descent.}
\end{theorem}

\begin{definition}[Derived completion]\label{def:derived-completion}
For $K\in D(X,A)$, regard $\operatorname{Red}_A(K)$ as an $A_X$-tower by
restriction of scalars and define
\[
 \widehat K^{\,L}:=
 \Rlim_{A_X}\bigl(\mathbf D\operatorname{Res}_A
       (\operatorname{Red}_A(K))\bigr)\in D(X,A).
\]
The quotient maps $K\to\operatorname{red}_n(K)$ are compatible and define
the completion unit
\[
 \eta_K:K\longrightarrow\widehat K^{\,L}.
\]
The object $K$ is \emph{derived $\varpi$-complete} if $\eta_K$ is an
isomorphism.  Let $D_{\mathrm{comp}}(X,A)$ be the full subcategory of such
objects.
\end{definition}

\begin{proposition}[Formal properties of completion]\label{prop:completion-formal}
The assignment $K\mapsto\widehat K^{\,L}$ is exact triangulated and
$\eta:\operatorname{id}\Rightarrow\widehat{(-)}^{\,L}$ is natural.  The
subcategory $D_{\mathrm{comp}}(X,A)$ is closed under shifts, cones, and
direct summands.
\depends{Proposition~\ref{prop:reductions-normalized};
Theorem~\ref{thm:milnor-triangle}; exactness and package independence of
$\Rlim$.}
\end{proposition}

\begin{definition}[Reconstruction from a normalized system]
\label{def:reconstruction}
For $K_\bullet\in\Dnorm(X,A_\bullet)$, define
\[
 \operatorname{Real}_A(K_\bullet)=
 \Rlim_{A_X}\bigl(\mathbf D\operatorname{Res}_A(K_\bullet)\bigr)
 \in D(X,A).
\]
Let
\[
 p_{K_\bullet,n}:\operatorname{Real}_A(K_\bullet)
 \longrightarrow \operatorname{Res}_{A/A_n}(K_n)
\]
be the derived-limit projection in $D(X,A)$.  Let
\[
 \mu_{K_n}:(A_n)_X\otimes^L_{A_X}
 \operatorname{Res}_{A/A_n}(K_n)\longrightarrow K_n
\]
be the derived $(A_n)_X$-action map.  Define the reconstruction morphism to
be the fully typed composite
\[
\begin{aligned}
 \varepsilon_{K_\bullet,n}:\quad
 \operatorname{red}_n(\operatorname{Real}_A(K_\bullet))
 &=(A_n)_X\otimes^L_{A_X}\operatorname{Real}_A(K_\bullet)\\
 &\xrightarrow{\operatorname{id}\otimes p_{K_\bullet,n}}
 (A_n)_X\otimes^L_{A_X}\operatorname{Res}_{A/A_n}(K_n)\\
 &\xrightarrow{\mu_{K_n}}K_n
 \qquad\text{in }D(X,A_n).
\end{aligned}
\]
For $K\in D_{\mathrm{comp}}(X,A)$, the completion unit is the natural map
\[
 \eta_K:K\longrightarrow
 \operatorname{Real}_A(\operatorname{Red}_A(K)).
\]
\end{definition}

\begin{theorem}[Reduction commutes with the relevant derived limit]
\label{thm:reduction-limit}
For every tower $C_\bullet\in D(\mathsf{Tow}_A(X))$ and every $m\geq1$, the
canonical exchange morphism
\[
 \operatorname{red}_m\!\left(\Rlim_{A_X}C_\bullet\right)
 \longrightarrow
 \Rlim_{(A_m)_X}\bigl(\mathsf T_mC_\bullet\bigr)
\]
is an isomorphism in $D(X,A_m)$.
\depends{Proposition~\ref{prop:quotient-resolution};
Theorems~\ref{thm:tower-reduction}, \ref{thm:milnor-triangle}, and
\ref{thm:perfect-tensor-limit}; perfection of the two-term $A_X$-complex
representing $(A_m)_X$.}
\end{theorem}

\begin{proposition}[Fixed-level change-of-rings decomposition]
\label{prop:fixed-level-change-rings}
Let $K_\bullet$ be normalized and fix $m\geq1$.  For $n\geq m$, the
two-term $A_X$-resolution of $(A_n)_X$ and normalization give an isomorphism
\[
 \operatorname{ev}_n\!\left(
   \mathsf T_m\mathbf D\operatorname{Res}_A(K_\bullet)\right)
 \xrightarrow{\sim}K_m\oplus K_m[1]
\]
in $D(X,A_m)$.  Under these isomorphisms, the transition from $n+1$ to $n$
is the identity on the summand $K_m$ and multiplication by $\varpi$ on the
summand $K_m[1]$.
\depends{Proposition~\ref{prop:quotient-resolution};
Definition~\ref{def:normalized-system}; the chain map
$L_{n+1}\to L_n$ whose degree $0$ component is the identity and whose
degree $-1$ component is multiplication by $\varpi$.}
\end{proposition}

\begin{lemma}[Derived limit of a nilpotent-transition tower]
\label{lem:nilpotent-tower}
Let $D\in D(X,A_m)$ and consider the constant-object tower
\[
 \cdots\xrightarrow{\varpi}D\xrightarrow{\varpi}D
 \xrightarrow{\varpi}D.
\]
Regarded as an object $N_\bullet\in
D(\mathsf{Tow}_{(A_m)_X}(X))$, it satisfies
\[
 \Rlim_{(A_m)_X}N_\bullet=0.
\]
More generally, the same holds for any object of
$D(\mathsf{Tow}_{(A_m)_X}(X))$ whose transition operator on the derived
product is uniformly nilpotent.
\depends{Theorem~\ref{thm:milnor-triangle}; the equality
$\varpi^m=0$ in $A_m$.}
\end{lemma}

\begin{theorem}[Normalized tail calculation]\label{thm:normalized-tail}
For a normalized system $K_\bullet$ and fixed $m$, there is a canonical
isomorphism
\[
 \Rlim_{(A_m)_X}\!\left(
 \mathsf T_m\mathbf D\operatorname{Res}_A(K_\bullet)
 \right)\xrightarrow{\sim}K_m.
\]
\depends{Proposition~\ref{prop:fixed-level-change-rings};
Lemma~\ref{lem:nilpotent-tower}; cofinal invariance of $\Rlim$.}
\end{theorem}

\begin{theorem}[Reconstruction counit]\label{thm:reconstruction-counit}
For every normalized system, the morphism
\[
 \varepsilon_{K_\bullet,n}:
 \operatorname{red}_n(\operatorname{Real}_A(K_\bullet))
 \xrightarrow{\sim}K_n
\]
is an isomorphism for every $n$, and these isomorphisms form an isomorphism
of normalized systems
\[
 \operatorname{Red}_A\operatorname{Real}_A(K_\bullet)
       \xrightarrow{\sim}K_\bullet.
\]
\depends{Definition~\ref{def:reconstruction};
Theorems~\ref{thm:reduction-limit} and \ref{thm:normalized-tail}.}
\end{theorem}

\begin{theorem}[Completeness of reconstruction]\label{thm:real-complete}
For every normalized system,
$\operatorname{Real}_A(K_\bullet)$ is derived $\varpi$-complete.  Its
completion unit is identified, after every finite reduction, with the
reconstruction counit of Theorem~\ref{thm:reconstruction-counit}.
\depends{Proposition~\ref{prop:completion-formal};
Theorem~\ref{thm:reconstruction-counit}; Definition~\ref{def:derived-completion}.}
\end{theorem}

\begin{theorem}[Idempotence of derived completion]
\label{thm:completion-idempotent}
For every $K\in D(X,A)$, its derived completion is derived complete, and
the unit
\[
 \eta_{\widehat K^{\,L}}:\widehat K^{\,L}
 \xrightarrow{\sim}\widehat{\widehat K^{\,L}}^{\,L}
\]
is an isomorphism.  The two natural maps from
$\widehat K^{\,L}$ to its double completion agree.
\depends{Proposition~\ref{prop:reductions-normalized};
Theorems~\ref{thm:reconstruction-counit} and \ref{thm:real-complete}.}
\end{theorem}

\begin{theorem}[Completion--system equivalence]\label{thm:adic-equivalence}
Reduction and reconstruction are quasi-inverse exact equivalences
\[
 \operatorname{Red}_A:
 D_{\mathrm{comp}}(X,A)\xrightarrow{\sim}
 \Dnorm(X,A_\bullet),
 \qquad
 \operatorname{Real}_A\ \text{quasi-inverse}.
\]
The unit is $\eta_K$ and the counit is
$(\varepsilon_{K_\bullet,n})_n$.  Restricting both sides to uniformly
constructible systems gives the corresponding $\Dbc$ equivalence, and
restricting to uniformly finite-Tor systems gives the $\Dctf$ equivalence.
\depends{Theorems~\ref{thm:reconstruction-counit},
\ref{thm:real-complete}, and \ref{thm:completion-idempotent};
derived completeness in
Definition~\ref{def:derived-completion}.}
\end{theorem}

\begin{corollary}[Conservativity of finite reductions]
\label{cor:reduction-conservative}
Let $u:K\to L$ be any morphism in $D_{\mathrm{comp}}(X,A)$.  If
\[
 \operatorname{red}_n(u):\operatorname{red}_n(K)
       \xrightarrow{\sim}\operatorname{red}_n(L)
\]
is an isomorphism for every $n$, then $u$ is an isomorphism.
\depends{Theorem~\ref{thm:adic-equivalence}; componentwise detection of
isomorphisms in the normalized system category.}
\end{corollary}

\section{Integral lisse and constructible categories}

\begin{definition}[Integral coefficient categories]\label{def:integral-dbc}
Define $\Dbc(X,A)$ to be the full subcategory of $D(X,A)$ consisting of
derived $\varpi$-complete objects $K$ such that
$\operatorname{Red}_A(K)$ is uniformly constructible.  Define
$\Dctf(X,A)$ similarly using uniformly finite-Tor systems.  A complete
complex is \emph{lisse} if its reduction system is uniformly lisse.
\end{definition}

\begin{corollary}[Integral constructibility has finite Tor amplitude]
\label{cor:integral-dbc-dctf}
For the pinned DVR coefficients,
\[
 \Dbc(X,A)=\Dctf(X,A).
\]
In particular, every integral constructible object may be used as the
finite-Tor first variable in the internal-Hom comparison below.
\depends{Definition~\ref{def:integral-dbc};
Proposition~\ref{prop:uniform-tor-dvr}.}
\end{corollary}

\begin{definition}[Completed tensor product and its unit]
\label{def:completed-tensor}
For $K,L\in D_{\mathrm{comp}}(X,A)$ define
\[
 K\widehat\otimes_A^LL
 :=\widehat{K\otimes^L_{A_X}L}^{\,L},
 \qquad
 \mathbf1_A:=\widehat{A_X}^{\,L}.
\]
If $P$ is a finite projective $A$-module, let $\underline P_X$ be the
$A_X$-module sheaf associated with
\[
 U\longmapsto\operatorname{Cont}(|U|,P),
\]
where $P$ has its $\varpi$-adic topology.  For a bounded complex $P^\bullet$
of finite projective $A$-modules, $\underline P_X$ denotes the corresponding
termwise sheaf complex with its induced differential.
\end{definition}

\begin{theorem}[Completed symmetric monoidal structure]
\label{thm:completed-monoidal}
The bifunctor $\widehat\otimes_A^L$ and the object $\mathbf1_A$ make
$D_{\mathrm{comp}}(X,A)$ an ordinary symmetric monoidal triangulated
category.  The associator, left and right unitors, and symmetry are the
completion of the corresponding maps for $\otimes^L_{A_X}$; they satisfy
the pentagon, triangle, and hexagon equations.  For every $n\geq1$ there is
a coherent natural isomorphism
\[
 \operatorname{red}_n(K\widehat\otimes_A^LL)
 \xrightarrow{\sim}
 \operatorname{red}_n(K)\otimes^L_{(A_n)_X}
 \operatorname{red}_n(L),
\]
compatible with the tensor unit, associator, symmetry, and all transition
maps $A_n\to A_m$.
\depends{Theorems~\ref{thm:completion-idempotent},
\ref{thm:reduction-limit}, and \ref{thm:adic-equivalence}; symmetric
monoidal coherence of derived tensor product.}
\end{theorem}

\begin{definition}[Duality datum]
\label{def:duality-datum}
A duality datum for $K\in D_{\mathrm{comp}}(X,A)$ consists of an object
$K^\vee$ and morphisms
\[
 \operatorname{coev}:\mathbf1_A\longrightarrow
 K\widehat\otimes_A^LK^\vee,
 \qquad
 \operatorname{ev}:K^\vee\widehat\otimes_A^LK
 \longrightarrow\mathbf1_A
\]
such that the composites
\[
\begin{aligned}
 K&\xrightarrow{\operatorname{coev}\widehat\otimes\operatorname{id}}
 K\widehat\otimes K^\vee\widehat\otimes K
 \xrightarrow{\operatorname{id}\widehat\otimes\operatorname{ev}}K,\\
 K^\vee&\xrightarrow{\operatorname{id}\widehat\otimes\operatorname{coev}}
 K^\vee\widehat\otimes K\widehat\otimes K^\vee
 \xrightarrow{\operatorname{ev}\widehat\otimes\operatorname{id}}K^\vee
\end{aligned}
\]
are the identity after inserting the associators and unitors of
Theorem~\ref{thm:completed-monoidal}.  The object $K$ is
\emph{dualizable} if it admits such a datum.
\end{definition}

\begin{theorem}[Local perfect criterion]\label{thm:local-perfect}
Assume that $X$ is locally noetherian and that the selected pro-etale basis
admits local refinements by quasi-compact weakly contractible objects.  Let
$K\in D_{\mathrm{comp}}(X,A)$.  The following are equivalent.
\begin{enumerate}
\item $K$ is lisse in the sense of Definition~\ref{def:integral-dbc}.
\item Pro-etale locally on $X$, there is a bounded complex $P$ of finite
projective $A$-modules and an isomorphism
$K\simeq\underline P_X$ in $D(X,A)$.
\item Pro-etale locally, $K$ is dualizable in the sense of
Definition~\ref{def:duality-datum}.
\end{enumerate}
Under these conditions all reductions have one common cohomological and
Tor-amplitude bound.
\depends{Theorems~\ref{thm:adic-equivalence} and
\ref{thm:completed-monoidal}; finite-level perfect-complex criterion;
weakly contractible local calculation.}
\end{theorem}

\begin{theorem}[Integral triangulated closure]\label{thm:integral-closure}
Assume that $X$ is topologically noetherian and that finite-level
constructibility is available.  The categories $\Dbc(X,A)$ and
$\Dctf(X,A)$ are full triangulated subcategories of $D(X,A)$ and are closed
under direct summands.  Every reduction restricts to an exact functor
\[
 \operatorname{red}_n:\Dbc(X,A)\longrightarrow\Dbc(X,A_n),
\qquad
 \operatorname{red}_n:\Dctf(X,A)\longrightarrow\Dctf(X,A_n).
\]
These functors retain the coherent transition isomorphisms
$\alpha_{n,m}$.
\depends{Theorems~\ref{thm:uniform-closure},
\ref{thm:adic-equivalence}, and \ref{thm:local-perfect};
Definition~\ref{def:integral-dbc}.}
\end{theorem}

\section{Rational localization and Karoubi completion}

\begin{definition}[Fraction morphisms]\label{def:fraction-category}
Let $\mathcal C=\Dbc(X,A)$.  A fraction representative from $K$ to $L$ is
a pair $(u,r)$ with $u\in\Hom_{\mathcal C}(K,L)$ and $r\geq0$, written
$u/\varpi^r$.  Define
\[
 (u,r)\sim(v,s)
 \quad\Longleftrightarrow\quad
 \text{there is $t\geq0$ such that }
 \varpi^{t+s}u=\varpi^{t+r}v.
\]
Set
\[
 \Hom_{\mathcal C[1/\varpi]}(K,L)
 =\{(u,r)\}/\!\sim.
\]
Addition and composition are
\[
 \frac{u}{\varpi^r}+\frac{v}{\varpi^s}
   =\frac{\varpi^su+\varpi^rv}{\varpi^{r+s}},
 \qquad
 \frac{v}{\varpi^s}\circ\frac{u}{\varpi^r}
   =\frac{v\circ u}{\varpi^{r+s}},
\]
and the identity of $K$ is $\operatorname{id}_K/1$.
\end{definition}

\begin{proposition}[The localized category]\label{prop:fraction-category}
The relation $\sim$ is an equivalence relation.  Addition, identities, and
composition in Definition~\ref{def:fraction-category} are independent of
representatives and make $\mathcal C[1/\varpi]$ an $E$-linear additive
category.  Its Hom groups identify canonically with
\[
 \Hom_{\mathcal C[1/\varpi]}(K,L)
 \cong\Hom_{\mathcal C}(K,L)\otimes_AE.
\]
\depends{Definition~\ref{def:central-scalar}; centrality of multiplication
by $\varpi$.}
\end{proposition}

\begin{theorem}[Triangulated central localization]
\label{thm:triangulated-localization}
There is a unique triangulated structure on $\mathcal C[1/\varpi]$ for which
the localization functor
\[
 Q_A:\mathcal C\longrightarrow\mathcal C[1/\varpi],
 \qquad K\longmapsto K,\qquad u\longmapsto u/1,
\]
is exact.  A triangle is distinguished precisely when it is isomorphic to
the localization of a distinguished triangle in $\mathcal C$.
\depends{Proposition~\ref{prop:fraction-category};
Theorem~\ref{thm:integral-closure}; the central localization theorem for
triangulated categories.}
\end{theorem}

\begin{definition}[Karoubi triples]\label{def:karoubi}
The Karoubi envelope $\Kar(\mathcal C[1/\varpi])$ has objects $(K,e)$ with
\[
 K\in\mathcal C[1/\varpi],\qquad
 e\in\operatorname{End}(K),\qquad e^2=e.
\]
A morphism $u:(K,e)\to(L,d)$ is a morphism
$u\in\Hom_{\mathcal C[1/\varpi]}(K,L)$ satisfying
\[
 u=d\,u\,e.
\]
The identity of $(K,e)$ is $e$, and composition is inherited from the
localized category.
\end{definition}

\begin{theorem}[Triangulated Karoubi completion]\label{thm:karoubi-triangulated}
The category $\Kar(\mathcal C[1/\varpi])$ has a canonical
idempotent-complete triangulated structure.  Its distinguished triangles
are the triangles isomorphic to direct summands, in the category of
triangles, of distinguished triangles in $\mathcal C[1/\varpi]$.  The
canonical functor
\[
 \mathcal C[1/\varpi]\longrightarrow\Kar(\mathcal C[1/\varpi]),
 \qquad K\longmapsto(K,\operatorname{id}_K),
\]
is exact.
\depends{Theorem~\ref{thm:triangulated-localization};
Definition~\ref{def:karoubi}; the idempotent-completion theorem for
triangulated categories.}
\end{theorem}

\begin{definition}[Rational constructible category and rationalization]
\label{def:rational-category}
Define
\[
 \Dbc(X,E):=\Kar\bigl(\Dbc(X,A)[1/\varpi]\bigr).
\]
The exact rationalization functor is
\[
 \operatorname{rat}_{E/A}:\Dbc(X,A)\longrightarrow\Dbc(X,E),
 \qquad K\longmapsto(K,\operatorname{id}_K),
\]
and similarly for $\Dctf$.  For $A=\Zl$ and $E=\Ql$ this is the coefficient
map $\Dbc(X,\Zl)\to\Dbc(X,\Ql)$.
\end{definition}

\begin{theorem}[Independence of the uniformizer]
\label{thm:uniformizer-independence}
Let $\varpi'=u\varpi$ be another uniformizer of $A$, with $u\in A^\times$.
The identity on objects and the correspondence
\[
 \frac{f}{\varpi^n}\longmapsto\frac{u^nf}{(\varpi')^n}
\]
give mutually inverse exact equivalences
\[
 \Dbc(X,A)[1/\varpi]\xrightarrow{\sim}
 \Dbc(X,A)[1/\varpi'].
\]
After Karoubi completion, this gives a canonical exact $E$-linear
equivalence between the two definitions of $\Dbc(X,E)$, compatible with
rationalization.
\depends{Definition~\ref{def:fraction-category};
Theorems~\ref{thm:triangulated-localization} and
\ref{thm:karoubi-triangulated}.}
\end{theorem}

\begin{definition}[Literal rational sheaf and lattices]
\label{def:literal-rational}
Let $E_X^{\mathrm{lit}}=A_X[1/\varpi]$ be the localization of the sheaf of
rings $A_X$.  Let $D(X,E_X^{\mathrm{lit}})$ be the derived category of
ordinary $E_X^{\mathrm{lit}}$-module sheaves.  An $A$-lattice in
$L\in D(X,E_X^{\mathrm{lit}})$ is a pair $(K,\lambda)$ with
$K\in\Dbc(X,A)$ and an isomorphism
\[
 \lambda:K\otimes^L_{A_X}E_X^{\mathrm{lit}}
       \xrightarrow{\sim}L.
\]
Let $\Dbclat(X,E_X^{\mathrm{lit}})$ be the full subcategory of
objects admitting such a lattice.

Let $\Dbclit(X,E_X^{\mathrm{lit}})$ be the full subcategory of
$D(X,E_X^{\mathrm{lit}})$ consisting of bounded complexes whose cohomology
sheaves are constructible and locally constant of finite-dimensional
$E$-vector spaces on a finite constructible stratification, with bounded
monodromy on every connected component of every stratum.  Here bounded
monodromy means that the image of the continuous pro-etale monodromy
representation is contained in a compact subgroup of
$\operatorname{GL}_r(E)$; equivalently, it stabilizes an $A$-lattice in
$E^r$.
\end{definition}

\begin{proposition}[Integral objects have admissible literal realizations]
\label{prop:integral-literal-landing}
For every $K\in\Dbc(X,A)$, the object
\[
 K_E:=K\otimes^L_{A_X}E_X^{\mathrm{lit}}
\]
is bounded, constructible, and has bounded monodromy on every lisse
stratum.  Thus scalar extension defines an additive functor
\[
 \Dbc(X,A)\longrightarrow\Dbclit(X,E_X^{\mathrm{lit}})
\]
that commutes with shifts and sends integral distinguished triangles to
distinguished triangles in the ambient derived category.
\depends{Theorems~\ref{thm:integral-closure} and
\ref{thm:local-perfect}; finite stratification and exactness of derived
localization $A_X\to A_X[1/\varpi]$.}
\end{proposition}

\begin{theorem}[Lattices for lisse sheaves]\label{thm:lisse-lattice}
Assume that $X$ is quasi-compact and locally connected in the pro-etale
topology.  Let $\mathcal L$ be a lisse finite-rank
$E_X^{\mathrm{lit}}$-module sheaf whose monodromy image on each connected
component is bounded.  Then $\mathcal L$ has, on each connected component, a lisse
$A_X$-submodule sheaf $\mathcal L_A$ such that
\[
 \mathcal L_A\otimes_{A_X}E_X^{\mathrm{lit}}
 \xrightarrow{\sim}\mathcal L.
\]
Two such lisse lattices are commensurable: there is $N\geq0$ such that,
inside $\mathcal L$,
\[
 \varpi^N\mathcal L_A\subseteq\mathcal L'_A
 \subseteq\varpi^{-N}\mathcal L_A.
\]
\depends{Compactness of the image of the continuous monodromy
representation; existence and commensurability of invariant lattices in a
finite-dimensional nonarchimedean representation.}
\end{theorem}

\begin{definition}[Integral lattice gluing property]
\label{def:lattice-gluing}
The pair $(X,A)$ has the integral lattice gluing property if the following
two statements hold.
\begin{enumerate}
\item If $K_1,K_2\in\Dbc(X,A)$ and
\[
 \delta:(K_2)_E\longrightarrow(K_1)_E[1]
\]
is a morphism after scalar extension, then there are $N\geq0$ and
$\delta_A:K_2\to K_1[1]$ whose scalar extension is
$\varpi^N\delta$.  The rationalized cone of $\delta_A[-1]$ is isomorphic to
the middle term of the triangle classified by $\delta$.
\item If $j:U\hookrightarrow X$ is a quasi-compact open immersion with
closed complement $i:Z\hookrightarrow X$, the functors $j_!$ and $i_*$ send
integral constructible lattices to integral constructible lattices and
their scalar extensions agree with the corresponding literal functors.
\end{enumerate}
In particular, for a triangle
\[
 j_!L_U\longrightarrow L\longrightarrow i_*L_Z
 \xrightarrow{\delta}j_!L_U[1]
\]
and lattices $K_U,K_Z$, these conditions produce a morphism
\[
 \delta_A:i_*K_Z\longrightarrow j_!K_U[1]
\]
whose scalar extension identifies with $\varpi^N\delta$.  The cone of
$\delta_A[-1]$ belongs to $\Dbc(X,A)$ and its scalar extension is
isomorphic to $L$.
\end{definition}

\begin{proposition}[Finite-level internal-Hom normalization]
\label{prop:internal-hom-normalization}
For $K,L\in\Dbc(X,A)$ put
\[
 \mathcal H_n(K,L)=
 R\mathcal Hom_{(A_n)_X}
   (\operatorname{red}_nK,\operatorname{red}_nL)
 \quad\in D(X,A_n).
\]
For $m\leq n$, perfect base change gives a natural isomorphism
\[
 \mathcal H_n(K,L)\otimes^L_{(A_n)_X}(A_m)_X
 \xrightarrow{\sim}\mathcal H_m(K,L).
\]
These maps satisfy the identity and triple-reduction equations and define a
normalized system
\[
 \mathcal H_\bullet(K,L)\in\Dnorm(X,A_\bullet).
\]
\depends{Proposition~\ref{prop:uniform-tor-dvr}; perfect base change for
internal Hom; coherence of finite reduction.}
\end{proposition}

\begin{definition}[Completed constructible internal Hom]
\label{def:completed-internal-hom}
For $K,L\in\Dbc(X,A)$ define
\[
 R\mathcal Hom_A(K,L):=
 \operatorname{Real}_A\bigl(\mathcal H_\bullet(K,L)\bigr)
 =\Rlim_{A_X}\mathbf D\operatorname{Res}_A
   \bigl(\mathcal H_\bullet(K,L)\bigr).
\]
\end{definition}

\begin{theorem}[Constructible internal Hom and bounded cohomology]
\label{thm:constructible-internal-hom}
Assume that $X$ is quasi-compact, quasi-separated, noetherian of finite
Krull dimension, that $\ell$ is invertible on $X$, and that constructible
finite-coefficient cohomology has a uniform finite cohomological-dimension
bound.  For $K,L\in\Dbc(X,A)$, the completed internal Hom of
Definition~\ref{def:completed-internal-hom} is derived complete and
constructible, and
$R\Gamma(X,R\mathcal Hom_A(K,L))$ is bounded with finitely generated
$A$-module cohomology.
\depends{Proposition~\ref{prop:internal-hom-normalization};
Theorems~\ref{thm:adic-equivalence},
\ref{thm:finite-closure}, and \ref{thm:reduction-limit}; finite-level
constructibility of internal Hom and finite cohomological dimension.}
\end{theorem}

\begin{theorem}[Comparison with the ambient internal Hom]
\label{thm:ambient-internal-hom-comparison}
Let $K\in\Dctf(X,A)$ and $L\in\Dbc(X,A)$.  The tensor--internal-Hom
adjunction gives a natural comparison
\[
 \chi_{K,L}:R\mathcal Hom_A(K,L)
 \longrightarrow R\mathcal Hom_{A_X}^{\mathrm{amb}}(K,L)
 \quad\text{in }D(X,A).
\]
The reduction of $\chi_{K,L}$ modulo $\varpi^n$ is the perfect base-change
isomorphism
\[
 \mathcal H_n(K,L)
 \xrightarrow{\sim}
 R\mathcal Hom_{(A_n)_X}
  (\operatorname{red}_nK,\operatorname{red}_nL).
\]
The ambient internal Hom on the right is derived complete, and
$\chi_{K,L}$ is an isomorphism.
\depends{Proposition~\ref{prop:internal-hom-normalization};
Theorems~\ref{thm:local-perfect}, \ref{thm:completed-monoidal}, and
\ref{thm:adic-equivalence}; conservativity of finite reductions from
Corollary~\ref{cor:reduction-conservative}.}
\end{theorem}

\begin{corollary}[Internal Hom represents derived morphisms]
\label{cor:internal-hom-represents}
For $K\in\Dctf(X,A)$, $L\in\Dbc(X,A)$, and $r\in\mathbf Z$, there is a
natural isomorphism
\[
 \Hom_{D(X,A)}(K,L[r])
 \xrightarrow{\sim}
 H^rR\Gamma\bigl(X,R\mathcal Hom_A(K,L)\bigr).
\]
It is compatible with composition, shifts, and finite reduction.
\depends{Theorem~\ref{thm:ambient-internal-hom-comparison}; the
tensor--internal-Hom adjunction on the ringed pro-etale site.}
\end{corollary}

\begin{theorem}[Localization of constructible derived Hom]
\label{thm:derived-hom-localization}
Under the hypotheses of
Theorem~\ref{thm:constructible-internal-hom}, let
$K,L\in\Dbc(X,A)$ and put
\[
 K_E=K\otimes^L_{A_X}E_X^{\mathrm{lit}},
 \qquad
 L_E=L\otimes^L_{A_X}E_X^{\mathrm{lit}}.
\]
Then the canonical morphism
\[
 R\Gamma(X,R\mathcal Hom_A(K,L))\otimes_A^LE
 \longrightarrow
 R\Gamma\!\left(X,
 R\mathcal Hom_{E_X^{\mathrm{lit}}}(K_E,L_E)\right)
\]
is an isomorphism in $D(E)$, naturally in $K,L$ and compatibly with shifts
in either variable.
\depends{Theorem~\ref{thm:constructible-internal-hom}; exactness of
$A\to E$; finite generation of the cohomology modules; scalar compatibility
of constructible internal Hom.}
\end{theorem}

\begin{theorem}[Hom localization]\label{thm:hom-localization}
Assume that $X$ is quasi-compact, quasi-separated, noetherian of finite
Krull dimension, $\ell$ is invertible on $X$, and constructible cohomology
has finite cohomological dimension.  For $K,L\in\Dbc(X,A)$, scalar extension
induces an isomorphism
\[
 \Hom_{\Dbc(X,A)}(K,L)\otimes_AE
 \longrightarrow
 \Hom_{D(X,E_X^{\mathrm{lit}})}
 (K\otimes_A^LE_X^{\mathrm{lit}},L\otimes_A^LE_X^{\mathrm{lit}})
\]
and the same statement holds with $L[r]$ in place of $L$ for every
$r\in\mathbf Z$.
\depends{Theorem~\ref{thm:derived-hom-localization};
Corollary~\ref{cor:internal-hom-represents};
Corollary~\ref{cor:integral-dbc-dctf}.}
\end{theorem}

\begin{proposition}[Clearing denominators in extension classes]
\label{prop:clear-extension-denominators}
Let $K_1,K_2\in\Dbc(X,A)$ and let
\[
 \delta:(K_2)_E\longrightarrow(K_1)_E[r]
\]
be a morphism in $D(X,E_X^{\mathrm{lit}})$.  There are $N\geq0$ and
$\delta_A:K_2\to K_1[r]$ such that
\[
 (\delta_A)_E=\varpi^N\delta.
\]
The statement is functorial under restriction to a constructible locally
closed subscheme and is compatible with composition after increasing $N$.
\depends{Theorem~\ref{thm:hom-localization}; the fraction description of
$M\otimes_AE$ for a finitely generated $A$-module $M$.}
\end{proposition}

\begin{theorem}[Integral open--closed preservation]
\label{thm:integral-open-closed}
Let $j:U\hookrightarrow X$ be a quasi-compact open immersion and
$i:Z\hookrightarrow X$ its closed complement.  Finite-level extension by
zero and closed pushforward assemble through normalized systems and
reconstruction to exact functors
\[
 j_!:\Dbc(U,A)\longrightarrow\Dbc(X,A),
 \qquad
 i_*:\Dbc(Z,A)\longrightarrow\Dbc(X,A).
\]
They commute with every finite reduction and with literal scalar extension:
\[
 (j_!K)_E\xrightarrow{\sim}j_!(K_E),
 \qquad
 (i_*L)_E\xrightarrow{\sim}i_*(L_E).
\]
\depends{Finite-level constructibility for $j_!$ and $i_*$; the
normalization theorem; Theorems~\ref{thm:adic-equivalence} and
\ref{thm:reconstruction-counit}; finite-level coefficient base change.}
\end{theorem}

\begin{proposition}[One-stratum and Postnikov lattice assembly]
\label{prop:postnikov-lattice}
Let $S$ be a connected stratum.  Every bounded
$E_X^{\mathrm{lit}}$-complex on $S$ whose cohomology sheaves are lisse with
bounded monodromy admits an integral lattice.  A finite Postnikov tower may
be chosen whose successive connecting morphisms are scalar extensions of
integral morphisms after multiplication by powers of $\varpi$.
\depends{Theorem~\ref{thm:lisse-lattice};
Proposition~\ref{prop:clear-extension-denominators}; the standard
truncation triangles.}
\end{proposition}

\begin{theorem}[The integral lattice gluing property]
\label{thm:lattice-gluing}
Assume the hypotheses of Theorem~\ref{thm:hom-localization} and the
finite-level constructibility hypotheses of
Theorem~\ref{thm:integral-open-closed}.  Then $(X,A)$ has the integral
lattice gluing property of Definition~\ref{def:lattice-gluing}.
\depends{Proposition~\ref{prop:clear-extension-denominators};
Theorem~\ref{thm:integral-open-closed}; exactness of cones in
$\Dbc(X,A)$.}
\end{theorem}

\begin{theorem}[Lattices for bounded constructible complexes]
\label{thm:complex-lattice}
Assume that $X$ is noetherian of finite Krull dimension, $\ell$ is
invertible on $X$, and the hypotheses of
Theorem~\ref{thm:lattice-gluing} hold.  Every
$L\in\Dbclit(X,E_X^{\mathrm{lit}})$ admits an integral lattice
$K\in\Dbc(X,A)$ with $K_E\xrightarrow{\sim}L$.
\depends{Proposition~\ref{prop:postnikov-lattice};
Theorem~\ref{thm:lattice-gluing}; finite common refinements of
constructible stratifications and open--closed localization triangles.}
\end{theorem}

\begin{theorem}[Triangulated closure of objects with lattices]
\label{thm:lattice-triangulated-closure}
Under Theorem~\ref{thm:hom-localization},
$\Dbclat(X,E_X^{\mathrm{lit}})$ is closed under shifts, finite direct sums,
and cones of all morphisms between its objects.  Its Karoubi envelope is an
idempotent-complete triangulated category.  The scalar-extension functor
\[
 \Dbc(X,A)[1/\varpi]\longrightarrow
 \Kar\bigl(\Dbclat(X,E_X^{\mathrm{lit}})\bigr)
\]
is exact.
\depends{Theorem~\ref{thm:hom-localization};
Proposition~\ref{prop:integral-literal-landing}; the cone construction in
the ambient derived category; the triangulated Karoubi-envelope theorem.}
\end{theorem}

\begin{corollary}[Literal constructible closure]
\label{cor:literal-constructible-closure}
Under Theorem~\ref{thm:complex-lattice},
\[
 \Dbclit(X,E_X^{\mathrm{lit}})=
 \Dbclat(X,E_X^{\mathrm{lit}}),
\]
so the literal constructible category is triangulated and its Karoubi
envelope is idempotent-complete triangulated.
\depends{Theorems~\ref{thm:complex-lattice} and
\ref{thm:lattice-triangulated-closure}.}
\end{corollary}

\begin{theorem}[Essential image after idempotent completion]
\label{thm:rational-essential-image}
Under the hypotheses of Theorems~\ref{thm:complex-lattice} and
\ref{thm:hom-localization}, scalar extension
\[
 \Dbc(X,A)[1/\varpi]\longrightarrow
 \Dbclit(X,E_X^{\mathrm{lit}})
\]
is fully faithful, its essential image contains every object admitting a
lattice, and every literal constructible complex is a direct summand of an
object in that image.
\depends{Theorems~\ref{thm:complex-lattice} and
\ref{thm:hom-localization}; Corollary~\ref{cor:literal-constructible-closure};
commensurability of lattices.}
\end{theorem}

\begin{theorem}[Integral--literal rational comparison]
\label{thm:literal-rational-comparison}
Under the hypotheses of Theorem~\ref{thm:rational-essential-image}, scalar
extension induces an exact equivalence
\[
 \Dbc(X,E)\xrightarrow{\sim}
 \Kar\bigl(\Dbclit(X,E_X^{\mathrm{lit}})\bigr).
\]
\depends{Proposition~\ref{prop:fraction-category};
Theorem~\ref{thm:karoubi-triangulated};
Theorems~\ref{thm:hom-localization} and
\ref{thm:rational-essential-image};
Corollary~\ref{cor:literal-constructible-closure}.}
\end{theorem}

\section{Finite extensions and coherent transition functors}

\begin{definition}[The finite-extension index]
\label{def:finite-extension-index}
Let $\mathsf{FinExt}_\ell$ be the poset category of finite intermediate
fields $\Ql\subset E\subset\Qlb$ in the fixed algebraic closure, with one
arrow $E\to F$ exactly when $E\subset F$.  The compositum of two objects is
again finite over $\Ql$, so $\mathsf{FinExt}_\ell$ is filtered.  For
$E\subset F$, put
\[
 A=\mathcal O_E,\qquad B=\mathcal O_F,
\]
choose uniformizers $\varpi_A,\varpi_B$, and write
\[
 \varpi_A B=v_{F/E}\varpi_B^{e(F/E)}B,
 \qquad v_{F/E}\in B^\times.
\]
The inclusion $A\hookrightarrow B$ is finite free.
\end{definition}

\begin{definition}[Coefficient sheaves for a finite extension]
\label{def:extension-coefficient-sheaves}
Define on the selected pro-etale basis
\[
 B_X(U)=\operatorname{Cont}(|U|,B),
 \qquad
 (B_n^{(A)})_X(U)=
 \operatorname{Cont}(|U|,B/\varpi_A^nB),
\]
and define $(B_m^{(B)})_X$ similarly using $B/\varpi_B^mB$.  The ring map
$A\to B$ gives morphisms
\[
 A_X\longrightarrow B_X,
 \qquad
 (A_n)_X\longrightarrow(B_n^{(A)})_X.
\]
Write $D(X,B)=D(X,B_X)$ and
$D(X,B_n^{(A)})=D(X,(B_n^{(A)})_X)$, and use the analogous notation for
$B_m^{(B)}$.
\end{definition}

\begin{theorem}[Finite freeness of the extended coefficient sheaf]
\label{thm:extension-sheaf-finite-free}
The presheaves of
Definition~\ref{def:extension-coefficient-sheaves} are sheaves.  Choice of
an $A$-basis of $B$ induces isomorphisms
\[
 A_X\otimes_AB\xrightarrow{\sim}B_X,
 \qquad
 (A_n)_X\otimes_{A_n}(B/\varpi_A^nB)
 \xrightarrow{\sim}(B_n^{(A)})_X.
\]
Consequently $B_X$ is finite locally free, hence perfect, as an
$A_X$-module, and
\[
 B_X\otimes^L_{A_X}(A_n)_X
 \xrightarrow{\sim}(B_n^{(A)})_X.
\]
All these isomorphisms commute with quotient transitions.
\depends{Theorem~\ref{thm:continuous-functions-sheaf}; finite freeness of
$B/A$; compatibility of continuous functions with a finite free basis.}
\end{theorem}

\begin{definition}[Cofinal quotient towers]\label{def:cofinal-towers}
Define
\[
 A_n=A/\varpi_A^nA,
 \qquad B_n^{(A)}=B/\varpi_A^nB,
 \qquad B_m^{(B)}=B/\varpi_B^mB.
\]
The inequalities
\[
 \varpi_B^{e(F/E)n+c}B\subseteq\varpi_A^nB
 \subseteq\varpi_B^{e(F/E)n-c}B
\]
for a fixed $c$ and all sufficiently large $n$ give cofinal maps between
the two $B$-quotient towers.  Their inverse limits both recover $B$.
\end{definition}

\begin{definition}[Integral scalar extension]\label{def:integral-extension}
For every $K\in D(X,A)$, define the ambient integral scalar extension
\[
 \operatorname{Ext}^{\mathrm{int}}_{B/A}(K)
 :=K\otimes^L_{A_X}B_X
 \quad\in D(X,B),
\]
using adapted K-flat derived scalar extension.  No completion is inserted
in the definition; preservation of completeness is a theorem below.
\end{definition}

\begin{theorem}[Integral scalar-extension coherence]
\label{thm:integral-extension-coherence}
For $E\subset F\subset G$ with valuation rings $A\subset B\subset C$,
derived tensor associativity gives a specified exact natural isomorphism
\[
 a^{\mathrm{int}}_{G,F,E}:
 \operatorname{Ext}^{\mathrm{int}}_{C/B}
 \operatorname{Ext}^{\mathrm{int}}_{B/A}
 \xRightarrow{\sim}
 \operatorname{Ext}^{\mathrm{int}}_{C/A}.
\]
Its component at $K$ is
\[
 (K\otimes^L_{A_X}B_X)\otimes^L_{B_X}C_X
 \xrightarrow{\sim}K\otimes^L_{A_X}C_X.
\]
There is a specified unit isomorphism at $E=F$.  The two composites for
$E\subset F\subset G\subset H$ are equal, and both unit diagrams commute.
\depends{Theorem~\ref{thm:derived-scalar-extension}; the tensor unit,
associator, and pentagon coherence.}
\end{theorem}

\begin{theorem}[Finite-level scalar comparison]
\label{thm:finite-scalar-comparison}
For every $n\geq1$ there is a natural isomorphism in
$D(X,B_n^{(A)})$
\[
 \operatorname{Ext}^{\mathrm{int}}_{B/A}(K)
     \otimes^L_{B_X}(B_n^{(A)})_X
 \xrightarrow{\sim}
 \bigl(K\otimes^L_{A_X}(A_n)_X\bigr)
     \otimes^L_{(A_n)_X}(B_n^{(A)})_X.
\]
These isomorphisms commute with all quotient transitions in $n$ and with
the cofinal change from the $\varpi_A$-tower to the $\varpi_B$-tower.
\depends{Theorem~\ref{thm:extension-sheaf-finite-free};
Definition~\ref{def:integral-extension}; associativity of derived tensor;
finite freeness of $B$ over $A$.}
\end{theorem}

\begin{definition}[The extended finite-reduction tower]
\label{def:extended-reduction-tower}
For $K\in D(X,A)$, let $\mathsf E_{B/A}(K)$ be the object of
$D(\mathsf{Tow}_{B_X}(X))$ whose $n$th evaluation is the restriction to
$B_X$ of
\[
 \operatorname{red}^A_n(K)
 \otimes^L_{(A_n)_X}(B_n^{(A)})_X
 \quad\in D(X,B_n^{(A)}),
\]
with transitions induced by the quotient maps
$B_{n+1}^{(A)}\to B_n^{(A)}$ and the reduction transitions of $K$.
\end{definition}

\begin{theorem}[Finite scalar extension preserves completeness]
\label{thm:scalar-extension-complete}
If $K$ is derived $\varpi_A$-complete, then
$K\otimes^L_{A_X}B_X$ is derived $\varpi_B$-complete.  More precisely,
finite freeness of $B_X$ over $A_X$ gives
\[
 \left(\Rlim_{A_X}
   \mathbf D\operatorname{Res}_A\operatorname{Red}_A(K)\right)
   \otimes^L_{A_X}B_X
 \xrightarrow{\sim}
 \Rlim_{B_X}\mathsf E_{B/A}(K),
\]
and cofinality identifies $\mathsf E_{B/A}(K)$ with the $\varpi_B$-adic
reduction tower of $K\otimes^L_{A_X}B_X$.
\depends{Theorems~\ref{thm:perfect-tensor-limit},
\ref{thm:extension-sheaf-finite-free}, and
\ref{thm:finite-scalar-comparison}; Definitions~\ref{def:cofinal-towers}
and \ref{def:extended-reduction-tower}.}
\end{theorem}

\begin{theorem}[Integral extension preserves admissibility]
\label{thm:integral-extension-preserves}
Assume that $X$ is topologically noetherian and finite-level
constructibility is available.  Then
$\operatorname{Ext}^{\mathrm{int}}_{B/A}$ restricts to exact functors
\[
 \Dbc(X,A)\longrightarrow\Dbc(X,B),
 \qquad
 \Dctf(X,A)\longrightarrow\Dctf(X,B).
\]
It preserves a common cohomological bound, sends a local perfect
$A$-complex to its scalar extension over $B$, and sends a witnessing finite
stratification to the same stratification.
\depends{Theorems~\ref{thm:finite-scalar-comparison},
\ref{thm:scalar-extension-complete},
\ref{thm:local-perfect}, and \ref{thm:adic-equivalence}.}
\end{theorem}

\begin{definition}[Localized and rational finite-extension functors]
\label{def:rational-extension}
Integral scalar extension sends multiplication by $\varpi_A$ to an
invertible scalar after localization at $\varpi_B$.  It therefore defines
an exact localized functor
\[
 \operatorname{Ext}^{\mathrm{loc}}_{F/E}:
 \Dbc(X,A)[1/\varpi_A]\longrightarrow
 \Dbc(X,B)[1/\varpi_B].
\]
For a fraction morphism its value is
\[
 \operatorname{Ext}^{\mathrm{loc}}_{F/E}
 \left(\frac{f}{\varpi_A^n}\right)
 =\frac{v_{F/E}^{-n}\operatorname{Ext}^{\mathrm{int}}_{B/A}(f)}
        {\varpi_B^{e(F/E)n}}.
\]
Define
\[
 \operatorname{Ext}_{F/E}:=
 \Kar(\operatorname{Ext}^{\mathrm{loc}}_{F/E}):
 \Dbc(X,E)\longrightarrow\Dbc(X,F).
\]
For a Karoubi object $(K,e)$, where
$e\in\End_{\Dbc(X,A)[1/\varpi_A]}(K)$ is a fraction idempotent,
\[
 \operatorname{Ext}_{F/E}(K,e)=
 \bigl(\operatorname{Ext}^{\mathrm{int}}_{B/A}K,
       \operatorname{Ext}^{\mathrm{loc}}_{F/E}(e)\bigr).
\]
\end{definition}

\begin{theorem}[Identity, composition, and scalar coherence]
\label{thm:extension-coherence}
There are specified natural isomorphisms
\[
 \iota_E:\operatorname{Ext}_{E/E}\xRightarrow{\sim}
     \operatorname{id}_{\Dbc(X,E)},
\]
and, for $E\subset F\subset G$,
\[
 a_{G,F,E}:
 \operatorname{Ext}_{G/F}\operatorname{Ext}_{F/E}
 \xRightarrow{\sim}\operatorname{Ext}_{G/E}.
\]
For $E\subset F\subset G\subset H$, the two composites from
\[
 \operatorname{Ext}_{H/G}\operatorname{Ext}_{G/F}
       \operatorname{Ext}_{F/E}
 \quad\text{to}\quad\operatorname{Ext}_{H/E}
\]
are equal, and the left and right unit diagrams commute.  The components
are the localized and Karoubi extensions of
$a^{\mathrm{int}}_{G,F,E}$.
\depends{Theorems~\ref{thm:integral-extension-preserves} and
\ref{thm:integral-extension-coherence};
Definition~\ref{def:rational-extension}; functoriality of central
localization and Karoubi completion.}
\end{theorem}

\begin{corollary}[The complete coefficient chain]
For $A=\Zl$ and $E=\Ql$, the constructions give exact functors
\[
 \Dbc(X,\Zl)
 \xrightarrow{\operatorname{rat}_{\Ql/\Zl}}
 \Dbc(X,\Ql)
 \xrightarrow{\operatorname{Ext}_{F/\Ql}}
 \Dbc(X,F)
\]
for every finite $F/\Ql$.  Finite reduction, rationalization, and finite
scalar extension are related by the comparison isomorphisms of
Theorems~\ref{thm:reduction-coherence} and
\ref{thm:finite-scalar-comparison}.
\depends{Definitions~\ref{def:rational-category} and
\ref{def:rational-extension}; Theorem~\ref{thm:extension-coherence}.}
\end{corollary}

\section{Reference points}

The finite-torsion input and the etale--pro-etale comparison are the
classical constructibility theorems transported to the pro-etale site.  The
complete-coefficient construction uses the replete pro-etale formalism of
Bhatt--Scholze.  The K-flat, K-injective, derived-limit, localization, and
Karoubi statements are ordinary results about abelian and triangulated
categories; parallel formulations occur throughout the Stacks Project.

\end{document}
